\section*{ID8633: \$\$P(s) \textbackslash{}propto s\textasciicircum{}\{-\textbackslash{}tau\}\$\$: P(s) ∝ s\textasciicircum{}-τ}
\paragraph{Metadata.}
PiID: 1276178583829 | RTEID: RTE:P33523.8562:T4.9715 | UUID: 13f7521e-be3d-516d-830c-bac3a01460bf

\paragraph{Canonical Statement.}
\$\$P(s) \textbackslash{}propto s\textasciicircum{}\{-\textbackslash{}tau\}\$\$

\paragraph{Canonical Equation.}
\[
P(s) \propto s^{-\tau}
\]

\paragraph{Dictionary.}
\$\$P(s) \textbackslash{}propto s\textasciicircum{}\{-\textbackslash{}tau\}\$\$: P(s) ∝ s\textasciicircum{}-τ — P(s) ∝ s\textasciicircum{}-τ

\paragraph{Encyclopedia.}
\$\$P(s) \textbackslash{}propto s\textasciicircum{}\{-\textbackslash{}tau\}\$\$: P(s) ∝ s\textasciicircum{}-τ is a scaling / order-of-magnitude relation, written P(s) ∝ s\textasciicircum{}-τ. It expresses a scaling or proportionality, capturing how the quantity grows with its parameters. It is built from τ, P, s. Within the Trawin Zero Point registry it serves as one element of the framework's mathematical narrative.
